本部分讨论智能体的“大脑”，以及智能体如何利用它进行学习。

上一部分介绍了预测下一词元的基础大语言模型如何训练。智能体模型通常会使用更先进的推理模型和多模态模型；本书仍将这些模型统称为“大语言模型”。本部分将讨论这些更先进的模型，说明为什么需要它们，以及如何训练和使用它们。

本部分第一章介绍推理模型。早期大语言模型无法解决简单的推理问题，思维链概念的提出推动了这一领域的重大进展。我们将讨论数学应用题、思维链提示与上下文学习，并考察模型如何通过自我反思来改进行为。自我反思循环进一步促成了大语言模型智能体控制框架的发展。

下一章把视野扩展到语言之外：我们将看到，Transformer 架构同样适用于视觉与动作领域，并讨论各种世界模型。这些模型帮助智能体在与世界交互时更明智地采取行动。

本部分最后一章介绍元学习。智能体大语言模型中的多种学习方式都体现了这一范式。我们将讨论元学习是什么，以及为什么它是一种富有挑战的机器学习形式。
